\section{De vuelta al cierre de la Sesión 6}

La Sesión 6 cerró la anti-equivalencia entre anillos y esquemas afines. Quedó la frase por formalizar: un esquema es un espacio localmente anillado que admite una cubierta por abiertos afines. Dos preguntas organizan esta sesión: si cada abierto básico $D(f)\subseteq\Spec A$ vuelve a ser afín y cómo construir un esquema a partir de piezas con reglas de traslape. El \cref{teo:s07-lema-pegado} responde la segunda pregunta: permite construir esquemas especificando sus piezas y las reglas de traslape.

\section{Restricción de un espacio localmente anillado a un abierto}

\begin{lema}[Restricción a un abierto]
\label{lema:s07-restriccion-abierto}
Sea $(X,\mathcal O_X)$ un espacio anillado y $U\subseteq X$ abierto. Para abiertos $V\subseteq U$, definamos $\mathcal O_X|_U(V)=\mathcal O_X(V)$, con las mismas restricciones. Entonces $(U,\mathcal O_X|_U)$ es un espacio anillado y, para cada $x\in U$, $ (\mathcal O_X|_U)_x\cong\mathcal O_{X,x}$. En particular, la restricción de un espacio localmente anillado sigue siendo localmente anillada.
\end{lema}
\begin{proof}
Los axiomas de gavilla se heredan de $\mathcal O_X$. Los abiertos de $U$ que contienen a $x$ son cofinales en los abiertos de $X$ que contienen a $x$: a una vecindad $W$ de $x$ en $X$ se le asocia $W\cap U$. Por ello los colímites que definen ambos tallos son canónicamente isomorfos.
\end{proof}

\begin{defi}[Subespacio localmente anillado abierto]
\label{defi:s07-subespacio-localmente-anillado-abierto}
En la situación descrita en el \cref{lema:s07-restriccion-abierto}, $(U,\mathcal O_X|_U)$ es el subespacio localmente anillado abierto determinado por $U$. La inclusión $j:U\hookrightarrow X$ induce $j^\#:\mathcal O_X\to j_*(\mathcal O_X|_U)$; sobre un abierto $V\subseteq X$, el mapa es la restricción $\mathcal O_X(V)\to\mathcal O_X(V\cap U)$. Es un morfismo de espacios localmente anillados cuando $X$ lo es.
\end{defi}

\section{Esquemas}

\begin{defi}[Esquema]
\label{defi:s07-esquema}
Un esquema es un espacio localmente anillado $(X,\mathcal O_X)$ en el que cada $x$ tiene un abierto $U\ni x$ isomorfo a $(\Spec A,\mathcal O_{\Spec A})$ como espacio localmente anillado, para algún anillo conmutativo $A$. Sus morfismos son los de espacios localmente anillados.
\end{defi}

\begin{obs}
\label{obs:s07-no-hace-falta-definicion-adicional}
No hace falta una definición adicional de morfismo de esquemas: es el morfismo de espacios localmente anillados entre los esquemas. Todo esquema afín es un esquema, tomando el abierto afín igual al espacio entero.
\end{obs}

\begin{ejem}[El espacio afín $\mathbb A_A^n$]
\label{ejem:s07-espacio-afin-n}
Para $A$ un anillo conmutativo y $n\geq0$, $\mathbb A_A^n=\Spec A[x_1,\ldots,x_n]$ ya es un esquema afín. Si $A=k$ es algebraicamente cerrado, sus puntos cerrados son los ideales $(x_1-a_1,\ldots,x_n-a_n)$ y se identifican con $k^n$. Para $A=\mathbb Z$, el espacio registra puntos sobre cada primo de $\Spec\mathbb Z$.
\end{ejem}

\section{Abiertos afines}

\begin{prop}[$D(f)$ es un abierto afín]
\label{prop:s07-d-f-abierto-afin}
Para $f\in A$, el abierto $D(f)\subseteq\Spec A$, con la gavilla estructural restringida, es isomorfo como espacio localmente anillado a $\Spec A_f$.
\end{prop}
\begin{proof}
El homeomorfismo se obtiene en el \cref{teo:s03-correspondencia-localizaciones} $\Spec A_f\cong D(f)$. En la base común formada por los abiertos $D(fg)$, la gavilla restringida de $\Spec A$ tiene secciones $A_{fg}$, mientras que la gavilla de $\Spec A_f$ tiene secciones $(A_f)_{\bar g}$. Ambos anillos son la $A$-álgebra universal en la que $f$ y $g$ son unidades, de modo que son canónicamente isomorfos. La unicidad del \cref{teo:s05-toda-gavilla-sobre-base-extiende} da el isomorfismo de gavillas y, por consiguiente, de espacios localmente anillados.
\end{proof}

\begin{prop}[Todo abierto de un esquema es un esquema]
\label{prop:s07-todo-abierto-esquema-esquema}
Todo abierto $V$ de un esquema $X$ es un esquema con estructura $\mathcal O_X|_V$.
\end{prop}
\begin{proof}
Para $x\in V$, elige un abierto afín $U\ni x$ en $X$. En $U\cap V$, que es abierto en $U\cong\Spec A$, existe un básico $D(g)$ que contiene a $x$. Por la \cref{prop:s07-d-f-abierto-afin}, $D(g)$ es afín. Estos abiertos dan una cubierta afín de $V$.
\end{proof}

\begin{defi}[Subesquema abierto, abierto afín, inmersión abierta]
\label{defi:s07-subesquema-abierto-abierto-afin-inmersion}
Si $V\subseteq X$ es abierto, $(V,\mathcal O_X|_V)$ es el subesquema abierto de $X$ determinado por $V$, y es un abierto afín si además es afín. Un morfismo de esquemas $Y\to X$ es una inmersión abierta si factoriza como un isomorfismo $Y\xrightarrow{\sim}V$ seguido de la inclusión de un subesquema abierto $V\subseteq X$.
\end{defi}

\begin{obs}
\label{obs:s07-abiertos-d-f-forman-base}
El homomorfismo $A\to A_f$ induce, por la \cref{prop:s07-d-f-abierto-afin}, el morfismo
\[
\Spec A_f\longrightarrow\Spec A
\]
es una inmersión abierta con imagen $D(f)$. Los abiertos $D(f)$ forman una base de abiertos afines en cada esquema afín. En cualquier esquema, cada punto tiene abiertos afines arbitrariamente pequeños: dentro de una carta afín se elige un básico contenido en el abierto dado. Por tanto los abiertos afines forman una base de la topología. El Lema de Comunicación Afín (Vakil, \S5.3), que permite verificar una propiedad en una sola cubierta afín, se enunciará más adelante.
\end{obs}

\begin{ejem}[La hipérbola]
\label{ejem:s07-hiperbola}
Sea $k$ un campo y $B=k[x,y]/(xy-1)$. El mapa $k[x]\to B$ induce una inmersión abierta $\Spec B\to\mathbb A_k^1$ con imagen $D(x)$, pues $B\cong k[x]_x$ y $y$ corresponde a $x^{-1}$.
\end{ejem}
\begin{obs}
\label{obs:s07-mismo-esquema-cerrado-k-2}
El mismo esquema es cerrado en $\mathbb A_k^2$, como conjunto de ceros de $xy-1$, y abierto en $\mathbb A_k^1$, como complemento del origen. Ser abierto o cerrado depende del esquema ambiente.
\end{obs}

\section{Datos de pegado}

Pegar objetos es construir un colímite. Para espacios topológicos $A,X,Y$ con mapas $A\to X$ y $A\to Y$, el pushout es el cociente $(X\amalg Y)/\sim$ con topología cociente, que identifica las dos imágenes de cada punto de $A$. Mapas continuos desde ese cociente equivalen a mapas desde $X$ y $Y$ compatibles sobre $A$.

\begin{ejem}[El círculo]
\label{ejem:s07-circulo}
Sean $I_1,I_2$ dos copias de $[0,1]$ y $A=\set{0,1}$, incluidas como extremos. El pushout $I_1\amalg_A I_2$ es homeomorfo a $S^1$: identifica los dos ceros y los dos unos, y los intervalos son los dos arcos.
\end{ejem}

En esquemas se pegan piezas por subesquemas abiertos y los cambios de carta son isomorfismos. La receta debe especificar qué partes se identifican y cómo. Por ejemplo, dos copias de $\mathbb A_k^1$ pueden pegarse a lo largo de $D(x)$ mediante la identidad o mediante $x\mapsto1/x$; son datos distintos y producen esquemas distintos.

\begin{defi}[Datos de pegado]
\label{defi:s07-datos-pegado}
Los datos de pegado para esquemas $\set{X_i}_{i\in I}$ consisten en subesquemas abiertos $X_{ij}\subseteq X_i$, con $X_{ii}=X_i$, e isomorfismos $\varphi_{ij}:X_{ij}\xrightarrow{\sim}X_{ji}$ que satisfacen:
\begin{enumerate}[label=(P\arabic*), leftmargin=2.1em]
\item $\varphi_{ii}=\id_{X_i}$ para todo $i$;
\item $\varphi_{ji}=\varphi_{ij}^{-1}$ para todo $i,j$;
\item para toda terna $i,j,k$, $\varphi_{ij}(X_{ij}\cap X_{ik})=X_{ji}\cap X_{jk}$ y, en $X_{ij}\cap X_{ik}$, $\varphi_{ik}=\varphi_{jk}\circ\varphi_{ij}$.
\end{enumerate}
\end{defi}
\begin{obs}
\label{obs:s07-primera-parte-p3-garantiza-composicion}
La primera parte de (P3) garantiza que la composición $\varphi_{jk}\circ\varphi_{ij}$ está definida en el traslape; la segunda dice que los dos cambios de carta coinciden.
\end{obs}
\begin{obs}[Con dos piezas no hay cociclo que verificar]
\label{obs:s07-dos-piezas-no-cociclo-verificar}
Si $I=\set{1,2}$, basta elegir $X_{12}\subseteq X_1$, $X_{21}\subseteq X_2$ y un isomorfismo $\varphi_{12}:X_{12}\to X_{21}$. Las condiciones diagonales e inversas determinan $\varphi_{11}$, $\varphi_{22}$ y $\varphi_{21}$; la condición (P3) se sigue de ellas.
\end{obs}

\section{El teorema central: el Lema de Pegado}

\begin{teo}[Lema de Pegado]
\label{teo:s07-lema-pegado}
Dados datos de pegado como en la \cref{defi:s07-datos-pegado}, existe un esquema $X$ y unas inmersiones abiertas $\psi_i:X_i\to X$ tales que
\begin{enumerate}[label=(\alph*)]
\item $X=\bigcup_i\psi_i(X_i)$;
\item $\psi_i(X_i)\cap\psi_j(X_j)=\psi_i(X_{ij})$ para todo $i,j$;
\item $\psi_j\circ\varphi_{ij}=\psi_i$ en $X_{ij}$.
\end{enumerate}
Además, si $(X',\psi_i')$ es otro esquema con inmersiones abiertas que satisfacen las mismas propiedades, existe un único isomorfismo $h:X\to X'$ con $h\circ\psi_i=\psi_i'$ para todo $i$.
\end{teo}
La demostración del \cref{teo:s07-lema-pegado} se desarrollará a continuación, en los cinco pasos de los \cref{sec:s07-paso-1,sec:s07-paso-2,sec:s07-paso-3,sec:s07-paso-4,sec:s07-paso-5}:
\begin{enumerate}
\item Construir el conjunto $X$ y las aplicaciones $\psi_i:X_i\to X$.
\item Dar a $X$ la topología cociente y probar que cada $\psi_i$ es un homeomorfismo sobre un abierto.
\item Construir $\mathcal O_X$ con familias compatibles de secciones y verificar los axiomas de gavilla.
\item Probar que $(X,\mathcal O_X)$ es un esquema y que las $\psi_i$ son inmersiones abiertas.
\item Probar la unicidad del pegado y del isomorfismo compatible con las $\psi_i$.
\end{enumerate}
\emph{Referencias clásicas:} Hartshorne, II.2, ejercicio 2.12; EGA I, \S2.3.

Como ejemplo guía, podemos pensar en dos copias de $\mathbb A_k^1$ pegadas por la identidad en $D(x)$. El conjunto resultante tendrá dos orígenes y una sola copia de cada punto de $D(x)$. El \cref{ejem:s07-recta-origen-doble} retomará esta construcción y probará que no es afín.

\subsection{Paso 1: el conjunto subyacente}
\label{sec:s07-paso-1}

Tomemos $X_0=\coprod_iX_i$ y definamos, para $x\in X_i$, $y\in X_j$,
\[
x\sim y\qquad\Longleftrightarrow\qquad x\in X_{ij}\quad\text{y}\quad y=\varphi_{ij}(x).
\]
La condición diagonal da reflexividad y la inversa $\varphi_{ji}=\varphi_{ij}^{-1}$ da simetría. Para transitividad, supongamos $x\sim y$ y $y\sim z$, con $x\in X_i$, $y\in X_j$, $z\in X_k$. Entonces $y=\varphi_{ij}(x)$ pertenece a $X_{ji}\cap X_{jk}$. Por la condición de traslape, existe $x'\in X_{ij}\cap X_{ik}$ con $y=\varphi_{ij}(x')$. La inyectividad de $\varphi_{ij}$ implica $x=x'$, así que $x\in X_{ik}$. La condición de cociclo da $\varphi_{ik}(x)=\varphi_{jk}(\varphi_{ij}(x))=z$, y por tanto $x\sim z$.

Definimos $X=X_0/\sim$ y $\psi_i$ como la composición de la inclusión de $X_i$ con la proyección al cociente. Cada $\psi_i$ es inyectiva: si $x,x'\in X_i$ son equivalentes, entonces $x'=\varphi_{ii}(x)=x$. Se tiene $\psi_j\circ\varphi_{ij}=\psi_i$. Toda clase tiene representante en alguna pieza, de modo que se cumple (a). Para (b), la contención $\psi_i(X_{ij})\subseteq\psi_i(X_i)\cap\psi_j(X_j)$ se sigue de (c). Recíprocamente, si $\psi_i(x)=\psi_j(y)$, entonces $x\sim y$, así que $x\in X_{ij}$. Esto prueba (a)--(c) para los conjuntos subyacentes.

\subsection{Paso 2: la topología}
\label{sec:s07-paso-2}

La topología de $X$ es la cociente de $q:X_0\to X$, la más fina que hace continua a $q$. Así, $W\subseteq X$ es abierto si y solo si $\psi_i^{-1}(W)$ es abierto en cada $X_i$. Para $V\subseteq X_i$ abierto,
\[
\psi_j^{-1}(\psi_i(V))=\varphi_{ij}(V\cap X_{ij}).
\]
Un punto $y\in X_j$ está en la preimagen de la izquierda si y solo si $y\sim x$ para algún $x\in V$. Esto equivale a que $y\in X_{ji}$ y $\varphi_{ji}(y)\in V$, o bien $y\in\varphi_{ij}(V\cap X_{ij})$.

Como $\varphi_{ij}$ es un isomorfismo entre abiertos, el lado derecho es abierto en $X_j$. La topología cociente implica que $\psi_i(V)$ es abierto en $X$. Por tanto $\psi_i$ es inyectiva, continua y abierta: un homeomorfismo sobre $\psi_i(X_i)$.

\subsection{Paso 3: la gavilla estructural}
\label{sec:s07-paso-3}

Para un isomorfismo $\varphi:Y\to Z$ y un abierto $U'\subseteq Z$, el transporte de secciones es
\[
\varphi^\sharp:\mathcal O_Z(U')\xrightarrow{\sim}\mathcal O_Y(\varphi^{-1}(U'))
\]
compatible con restricciones. El transporte lleva secciones de $Z$ a secciones de $Y$; esta dirección determina la condición de compatibilidad en los traslapes.

\begin{defi}[Imagen directa]
\label{defi:s07-imagen-directa}
Si $\psi:Y\to X$ es continua y $\mathcal F$ es una gavilla en $Y$, su imagen directa es la gavilla $\psi_*\mathcal F$ en $X$ dada por $(\psi_*\mathcal F)(W)=\mathcal F(\psi^{-1}(W))$, con las restricciones inducidas por las de $\mathcal F$.
\end{defi}

\begin{obs}
\label{obs:s07-imagen-directa-pushout}
La gavilla $(\psi_i)_*\mathcal O_{X_i}$ se llama imagen directa o \emph{pushforward}. El pushout topológico de la construcción anterior es un colímite que identifica puntos. Las secciones de la \cref{defi:s07-gavilla-estructural-pegado} se describen como familias compatibles en el producto
\[
\prod_i\bigl((\psi_i)_*\mathcal O_{X_i}\bigr)(W)
\]
\end{obs}

\begin{defi}[La gavilla estructural del pegado]
\label{defi:s07-gavilla-estructural-pegado}
Para un abierto $W\subseteq X$, definimos
\[
\mathcal O_X(W):=\set*{(s_i)_{i\in I}\in\prod_{i\in I}\mathcal O_{X_i}(\psi_i^{-1}(W))\given
s_i|_{\psi_i^{-1}(W)\cap X_{ij}}=\varphi_{ij}^{\sharp}\bigl(s_j|_{\psi_j^{-1}(W)\cap X_{ji}}\bigr)\text{ para todo }i,j}.
\]
\end{defi}

\begin{obs}
\label{obs:s07-condicion-gavilla-pegado}
Por la propiedad (c) del \cref{teo:s07-lema-pegado}, $\varphi_{ij}$ lleva $\psi_i^{-1}(W)\cap X_{ij}$ sobre $\psi_j^{-1}(W)\cap X_{ji}$. Así, ambos lados de la condición de la \cref{defi:s07-gavilla-estructural-pegado} son secciones sobre el mismo abierto. Además, $\mathcal O_X(W)$ es un subanillo del producto, pues las condiciones son igualdades entre homomorfismos de anillos. Falta verificar que sea una pregavilla y satisfaga los axiomas de gavilla.
\end{obs}

\begin{prop}
\label{prop:s07-gavilla-estructural}
La asignación $W\mapsto\mathcal O_X(W)$ de la \cref{defi:s07-gavilla-estructural-pegado}, con las restricciones coordenada a coordenada, es una gavilla de anillos en $X$.
\end{prop}
\begin{proof}
\emph{Pregavilla.} Si $W_0\subseteq W$, restringir cada coordenada lleva una familia compatible sobre $W$ a otra sobre $W_0$, pues los mapas $\varphi_{ij}^{\sharp}$ conmutan con las restricciones. La identidad y la composición se verifican coordenada a coordenada, porque son las de las gavillas $\mathcal O_{X_i}$. Por tanto, $\mathcal O_X$ es una pregavilla.

\emph{Identidad.} Sean dos secciones de $\mathcal O_X(W)$ que coinciden en cada abierto $W_\alpha$ de una cubierta de $W$. Fijemos $i$. Sus coordenadas $i$ coinciden en cada $\psi_i^{-1}(W_\alpha)$, que cubren $\psi_i^{-1}(W)$. El axioma de identidad para $\mathcal O_{X_i}$ implica que coinciden en todo $\psi_i^{-1}(W)$. Esto vale para cada $i$, así que las secciones son iguales.

\emph{Pegado.} Sean familias compatibles $(s_i^\alpha)_{i\in I}\in\mathcal O_X(W_\alpha)$ sobre una cubierta $W=\bigcup_\alpha W_\alpha$. Para cada $i$, el axioma de pegado de $\mathcal O_{X_i}$ produce una única sección $s_i\in\mathcal O_{X_i}(\psi_i^{-1}(W))$ cuyas restricciones son $s_i^\alpha$. Falta verificar la condición entre $s_i$ y $s_j$. En $\psi_i^{-1}(W)\cap X_{ij}$, los dos lados coinciden al restringirse a cada $\psi_i^{-1}(W_\alpha)\cap X_{ij}$, pues cada familia $(s_i^\alpha)_i$ ya cumplía esa condición. Estos abiertos cubren el traslape; el axioma de identidad para $\mathcal O_{X_i}$ da la igualdad en todo él. Así $(s_i)_i\in\mathcal O_X(W)$ y es la única sección que restringe a las familias dadas. Por tanto, $\mathcal O_X$ es una gavilla.
\end{proof}

\subsection{Paso 4: el pegado es un esquema}
\label{sec:s07-paso-4}

Fijemos $i\in I$ y un abierto $V\subseteq X_i$, y pongamos $W=\psi_i(V)$. Por el \cref{sec:s07-paso-2},
\[
V_j:=\psi_j^{-1}(W)=\varphi_{ij}(V\cap X_{ij})
\]
es abierto en $X_j$, y $V_i=V$.

\begin{lema}
\label{lema:s07-proyeccion-coordenada-iso}
La proyección
\[
\pi_i:\mathcal O_X(\psi_i(V))\longrightarrow\mathcal O_{X_i}(V),
\qquad (s_j)_{j\in I}\longmapsto s_i,
\]
es un isomorfismo de anillos.
\end{lema}
\begin{proof}
\emph{Inyectividad.} La condición de compatibilidad para el par $(i,j)$ determina $s_j$ a partir de $s_i$. Como $\varphi_{ij}$ es un isomorfismo, dos familias con igual coordenada $i$ son iguales.

\emph{Sobreyectividad.} Sea $s\in\mathcal O_{X_i}(V)$. La única elección posible, despejando la condición para el par $(i,j)$, es
\[
s_i=s,\qquad s_j=(\varphi_{ij}^{\sharp})^{-1}\bigl(s|_{V\cap X_{ij}}\bigr).
\]
La definición de $\mathcal O_X$ exige comprobar la compatibilidad para todos los pares.

\emph{Compatibilidad del par $(j,k)$.} Sean $T=V_j\cap X_{jk}$ y $T'=\varphi_{ji}(T)$, y pongamos $\sigma=s|_{T'}$. Por (P3), $T'\subseteq V\cap X_{ij}\cap X_{ik}$ y $\varphi_{jk}(T)=\varphi_{ik}(T')$. La definición de $s_j$ y $s_k$ da
\[
s_j|_T=\varphi_{ji}^{\sharp}(\sigma),\qquad
s_k|_{\varphi_{jk}(T)}=\varphi_{ki}^{\sharp}(\sigma).
\]
Por tanto,
\[
\varphi_{jk}^{\sharp}\bigl(s_k|_{\varphi_{jk}(T)}\bigr)
=(\varphi_{ki}\circ\varphi_{jk})^{\sharp}(\sigma)
=\varphi_{ji}^{\sharp}(\sigma)=s_j|_T.
\]
La segunda igualdad usa $\varphi_{jk}=\varphi_{ik}\circ\varphi_{ji}$ y $\varphi_{ki}=\varphi_{ik}^{-1}$. Así, la familia construida pertenece a $\mathcal O_X(W)$ y $\pi_i$ es sobreyectiva.
\end{proof}

\begin{prop}
\label{prop:s07-pegado-esquema}
El espacio $(X,\mathcal O_X)$ es un esquema, y cada $\psi_i:X_i\to X$ es una inmersión abierta que cumple las propiedades (a), (b) y (c) del \cref{teo:s07-lema-pegado}.
\end{prop}
\begin{proof}
\emph{De $\pi_i$ a $\psi_i^\sharp$.} Las proyecciones a la coordenada $i$ conmutan con las restricciones y definen, para cualquier abierto $W\subseteq X$, el mapa estructural
\[
\psi_i^\sharp(W):\mathcal O_X(W)\longrightarrow
\mathcal O_{X_i}(\psi_i^{-1}(W)),\qquad (s_j)_j\longmapsto s_i.
\]
Sobre $W=\psi_i(V)$, este mapa es precisamente $\pi_i$. El \cref{lema:s07-proyeccion-coordenada-iso} muestra que ese mapa es un isomorfismo; junto con el homeomorfismo sobre un abierto del \cref{sec:s07-paso-2}, prueba que $\psi_i$ es una inmersión abierta. Las propiedades (a), (b) y (c) ya valían para los conjuntos subyacentes, y la compatibilidad de las coordenadas prueba (c) como igualdad de morfismos: en cada traslape,
\[
(\psi_j\circ\varphi_{ij})^\sharp
=\varphi_{ij}^\sharp\circ\psi_j^\sharp=\psi_i^\sharp.
\]

\emph{$X$ es un esquema.} Para $x\in X_i$, el tallo $\mathcal O_{X,\psi_i(x)}$ es isomorfo a $\mathcal O_{X_i,x}$, que es local, pues $\psi_i(X_i)$ es abierto y $\psi_i$ es un isomorfismo sobre su imagen. Así, $X$ es localmente anillado. Los abiertos $\psi_i(X_i)$ cubren $X$. Como cada $X_i$ es un esquema y $\psi_i$ es un isomorfismo sobre su imagen, cada punto de $X$ tiene una base de abiertos afines. Por la \cref{defi:s07-esquema}, queda probada la existencia.
\end{proof}

\subsection{Paso 5: unicidad}
\label{sec:s07-paso-5}

\begin{prop}
\label{prop:s07-unicidad-pegado}
Si $(X',\psi_i')$ es otro esquema con inmersiones abiertas que cumplen las propiedades (a), (b) y (c) del \cref{teo:s07-lema-pegado}, existe un único isomorfismo $h:X\to X'$ tal que $h\circ\psi_i=\psi_i'$ para todo $i$.
\end{prop}
\begin{proof}
Definimos $h(\psi_i(x))=\psi_i'(x)$. Si $\psi_i(x)=\psi_j(y)$, la relación del primer paso dice que $x\in X_{ij}$ e $y=\varphi_{ij}(x)$. Por (c) en $X'$, $\psi_i'(x)=\psi_j'(y)$; así, la definición no depende del representante.

Para construir la inversa hay que probar el recíproco. Si $\psi_i'(x)=\psi_j'(y)$, la propiedad (b) en $X'$ da un $z\in X_{ij}$ con $\psi_i'(x)=\psi_i'(z)$. La inyectividad de $\psi_i'$ implica $x=z\in X_{ij}$. Por (c),
\[
\psi_j'(\varphi_{ij}(x))=\psi_i'(x)=\psi_j'(y),
\]
y la inyectividad de $\psi_j'$ da $y=\varphi_{ij}(x)$. Por tanto $x\sim y$ y queda bien definida $\bar h(\psi_i'(x))=\psi_i(x)$. Las identidades $h\circ\psi_i=\psi_i'$ y $\bar h\circ\psi_i'=\psi_i$ prueban que $h$ y $\bar h$ son inversas.

En cada abierto $\psi_i(X_i)$, el mapa es $\psi_i'\circ\psi_i^{-1}$, un isomorfismo de esquemas. Estos mapas coinciden en los traslapes por (c), así que se pegan a un isomorfismo $h:X\to X'$. La unicidad se sigue de que los $\psi_i(X_i)$ cubren $X$.
\end{proof}

\begin{proof}[Conclusión de la demostración]
Los \cref{sec:s07-paso-1,sec:s07-paso-2} construyen el espacio topológico $X$ y las aplicaciones $\psi_i$. La \cref{prop:s07-gavilla-estructural} verifica los axiomas de gavilla para la \cref{defi:s07-gavilla-estructural-pegado}; el \cref{lema:s07-proyeccion-coordenada-iso} y la \cref{prop:s07-pegado-esquema} muestran que $X$ es un esquema y que las $\psi_i$ son las inmersiones abiertas requeridas. La \cref{prop:s07-unicidad-pegado} da la unicidad. Esto demuestra el \cref{teo:s07-lema-pegado}.
\end{proof}

\section{El pegado como colímite}

\begin{prop}[El pegado es un coigualador]
\label{prop:s07-pegado-coigualador}
Sean $E=\coprod_{i,j}X_{ij}$ y $X_\amalg=\coprod_iX_i$, con mapas $s,t:E\rightrightarrows X_\amalg$ dados respectivamente por la inclusión de $X_{ij}$ en $X_i$ y por $\varphi_{ij}$ seguida de la inclusión en $X_j$. Los $\psi_i$ definen $\psi:X_\amalg\to X$, con $\psi\circ s=\psi\circ t$. El diagrama $E\rightrightarrows X_\amalg\xrightarrow{\psi}X$ es un coigualador en la categoría de esquemas. Es decir, para todo esquema $Y$ y todo $f:X_\amalg\to Y$ con $f\circ s=f\circ t$, existe un único $\bar f:X\to Y$ con $\bar f\circ\psi=f$.
\end{prop}
\begin{proof}
Un morfismo $X_\amalg\to Y$ equivale a una familia $f_i:X_i\to Y$. La igualdad $f\circ s=f\circ t$ dice exactamente que $f_i|_{X_{ij}}=f_j\circ\varphi_{ij}$. En los abiertos $\psi_i(X_i)$, los mapas $f_i\circ\psi_i^{-1}$ coinciden en los traslapes por esta igualdad, así que se pegan a un único morfismo $X\to Y$. La unicidad se sigue de que los $\psi_i(X_i)$ cubren $X$.
\end{proof}

\begin{obs}
\label{obs:s07-dar-morfismo-x-equivale-dar}
Dar un morfismo $X\to Y$ equivale a dar una familia $f_i:X_i\to Y$ con $f_i|_{X_{ij}}=f_j\circ\varphi_{ij}$. El argumento usa solo las propiedades (a), (b), (c), por lo que vale para cualquier pegado que las satisfaga y expresa de nuevo la unicidad.
\end{obs}

La distinción con los límites se mantiene: en la categoría de esquemas, los límites finitos existen por el objeto final $\Spec\mathbb Z$ junto con los productos fibrados. El objeto inicial es el esquema vacío.

\section{Consistencia: pegado y cubierta}

\begin{coro}
\label{coro:s07-x-esquema-cubierta-abiertos-u}
Sea $(X,\mathcal O_X)$ un esquema con una cubierta por abiertos $\set{U_i}_{i\in I}$. Tomando $U_{ij}=U_i\cap U_j$ y los mapas de transición identidad, los datos satisfacen la \cref{defi:s07-datos-pegado} y el esquema producido por el \cref{teo:s07-lema-pegado} es canónicamente isomorfo a $X$.
\end{coro}
\begin{proof}
El mapa de cada pieza $U_i$ a $X$ es una inmersión abierta. Como los $U_i$ cubren $X$ y sus intersecciones son exactamente $U_{ij}$, estos mapas satisfacen las tres propiedades del \cref{teo:s07-lema-pegado}. Su unicidad da el isomorfismo canónico.
\end{proof}
\begin{obs}
\label{obs:s07-definicion-esquema-cubiertas-afines-teorema}
En particular, el corolario se aplica a una cubierta por abiertos afines. La \cref{defi:s07-esquema} y el \cref{teo:s07-lema-pegado} describen la misma clase de objetos desde dos direcciones. Todo esquema resulta de pegar su propia cubierta afín; el pegado también permite construir esquemas cuando las piezas no están previamente dentro de uno común.
\end{obs}

\begin{samepage}
\section{Un ejemplo: la recta con origen doble}

\begin{ejem}[La recta con origen doble]
\label{ejem:s07-recta-origen-doble}
Sea $k$ un campo, $X_1=X_2=\mathbb A_k^1=\Spec k[x]$, $X_{12}=X_{21}=D(x)$ y $\varphi_{12}=\id$. Por la \cref{obs:s07-dos-piezas-no-cociclo-verificar} son datos de pegado. Sea $X$ el esquema que produce el \cref{teo:s07-lema-pegado}, con $\psi_1,\psi_2$. Los puntos $O_1:=\psi_1(0)$ y $O_2:=\psi_2(0)$ son distintos porque $0\notin X_{12}$. Fuera de los orígenes hay una sola copia de cada punto. Es la recta con origen doble, un esquema que no es afín.

\begin{center}
\begin{tikzpicture}[>=to, every node/.style={font=\small}]
  \draw[->] (-3,0)--(-0.25,0);
  \draw[->] (0.25,0)--(3,0);
  \draw[->] (-3,-0.45)--(-0.25,-0.45);
  \draw[->] (0.25,-0.45)--(3,-0.45);
  \fill[rojoCIMAT] (0,0) circle (2pt);
  \fill[grisCIMAT] (0,-0.45) circle (2pt);
  \node[above] at (0,0) {$O_1$};
  \node[below] at (0,-0.45) {$O_2$};
  \node[right] at (3,0) {$X_1$};
  \node[right] at (3,-0.45) {$X_2$};
  \node at (0.9,-1.05) {Las dos copias se identifican en $D(x)$};
\end{tikzpicture}
\end{center}

Por la \cref{defi:s07-gavilla-estructural-pegado}, una sección global es un par de polinomios en $k[x]$ que coinciden en $k[x]_x$; la inclusión $k[x]\hookrightarrow k[x]_x$ es inyectiva, así que los dos polinomios son iguales. Por tanto $\Gamma(X,\mathcal O_X)=k[x]$. Si $X$ fuera afín, sería isomorfo a $\Spec\Gamma(X,\mathcal O_X)=\Spec k[x]$. El mapa canónico al espectro de secciones globales viene dado por
\[
p\longmapsto\set*{s\in\Gamma(X,\mathcal O_X)\given s_p\in\mathfrak m_p},
\]
donde $\mathfrak m_p$ es el ideal maximal de $\mathcal O_{X,p}$. Ambos orígenes van al ideal $(x)$, de modo que este mapa no es inyectivo y $X$ no es afín.
\end{ejem}
\end{samepage}

\begin{obs}[El colímite entre esquemas afines]
\label{obs:s07-colimite-entre-esquemas-afines}
La categoría de esquemas afines es equivalente a la opuesta de la categoría de anillos. El pushout de $\Spec T\to\Spec R$ y $\Spec T\to\Spec S$ calculado entre afines es $\Spec(R\times_T S)$. Para $R=S=k[x]$, $T=k[x]_x$ y las inclusiones, $k[x]\times_{k[x]_x}k[x]\cong k[x]$. Por tanto, el pushout afín del diagrama
\[
\Spec k[x]\longleftarrow\Spec k[x]_x\longrightarrow\Spec k[x]
\]
es $\Spec k[x]$; en esquemas, el pushout es la recta con origen doble. Por la \cref{prop:s07-pegado-coigualador}, para dos piezas se obtiene la propiedad universal del pushout. La recta con origen doble se aplica sobre $\mathbb A_k^1$ sin ser isomorfa a ella. La inclusión de esquemas afines en esquemas no preserva este colímite. Por eso la demostración del \cref{teo:s07-lema-pegado} construye primero el espacio y la gavilla, y reconoce después los abiertos afines.
\end{obs}

\section{Ejercicios}

\begin{ejer}[Pegado con una sola pieza]
\label{ejer:s07-pegado-sola-pieza}
Verifica que si $I=\set{1}$ (un solo esquema $X_1$, con $X_{11}=X_1$ y $\varphi_{11}=\id$), el esquema producido por el \cref{teo:s07-lema-pegado} es canónicamente isomorfo a $X_1$, con $\psi_1$ la identidad.
\end{ejer}

\begin{samepage}
\begin{ejer}[Dos copias de $\Spec A$ pegadas por la identidad]
\label{ejer:s07-dos-copias-pegadas-identidad}
Sean $X_1=X_2=\Spec A$ dos copias, $f\in A$, $X_{12}=X_{21}=D(f)$ y el isomorfismo de pegado la identidad. Describe el esquema pegado en cada caso.
\begin{enumerate}
\item $f$ es invertible en $A$.
\item $A=k[x]$ con $k$ un campo y $f=x$. Pista: los puntos de $\Spec k[x]$ que no están en $D(x)$ no se identifican entre las dos copias.
\end{enumerate}
\end{ejer}
\end{samepage}

\begin{ejer}[El plano con recta doble]
\label{ejer:s07-plano-recta-doble}
Sea $k$ un campo, $X_1=X_2=\mathbb A_k^2=\Spec k[x,y]$, $X_{12}=X_{21}=D(x)$ y $\varphi_{12}=\id_{D(x)}$. Sea $X$ el esquema pegado.
\begin{enumerate}
\item Calcula $\Gamma(X,\mathcal O_X)$.
\item Muestra que $X$ no es afín, adaptando el argumento del \cref{ejem:s07-recta-origen-doble}. Pista: considera el ideal primo $(x)\subseteq k[x,y]$ visto en cada una de las dos copias.
\item Explica por qué esto no contradice el \cref{coro:s07-x-esquema-cubierta-abiertos-u}.
\end{enumerate}
\end{ejer}

\begin{ejer}[La imagen directa de una gavilla]
\label{ejer:s07-imagen-directa-gavilla}
Sean $\psi:Y\to X$ continua y $\mathcal F$ una gavilla en $Y$. Verifica que $\psi_*\mathcal F$ de la \cref{defi:s07-imagen-directa} es una gavilla y que la \cref{defi:s07-gavilla-estructural-pegado} da el subanillo de $\prod_i((\psi_i)_*\mathcal O_{X_i})(W)$ formado por las familias compatibles.
\end{ejer}

\begin{ejer}[Restricción de un morfismo a un abierto]
\label{ejer:s07-restriccion-morfismo-abierto}
Sea $f:X\to Y$ un morfismo de esquemas y $V\subseteq Y$ abierto. Verifica que $f^{-1}(V)$ es abierto y que la restricción induce un morfismo de esquemas $f|_{f^{-1}(V)}:f^{-1}(V)\to V$.
Usa la \cref{prop:s07-todo-abierto-esquema-esquema} para que $f^{-1}(V)$ y $V$ sean esquemas, y el \cref{lema:s07-restriccion-abierto} para sus estructuras locales. \emph{[Cf. Vakil, ejercicio de restricción de morfismos.]}
\end{ejer}
